% base_datos.tex — Capítulo 5: Base de Datos
\chapter{Base de Datos}

\section{Visión General}

La base de datos \textbf{finara\_db} está implementada en MySQL 8.0 con motor InnoDB. Consta de \textbf{10 tablas}, relaciones con integridad referencial mediante claves foráneas, 2 triggers automáticos y zona horaria configurada en la conexión del pool.

\begin{longtable}{>{\ttfamily}p{4.5cm} >{\centering}p{1.5cm} p{8cm}}
  \caption{Tablas de la base de datos finara\_db}
  \label{tab:tablas} \\
  \toprule
  \textbf{Tabla} & \textbf{Filas aprox.} & \textbf{Propósito} \\
  \midrule
  \endfirsthead
  \toprule
  \textbf{Tabla} & \textbf{Filas aprox.} & \textbf{Propósito} \\
  \midrule
  \endhead
  Usuarios                   & 1--100    & Cuentas de usuario con flags de onboarding \\
  Categorias                 & 10--50    & Categorías de gasto predefinidas y personalizadas \\
  Presupuestos               & 0--20     & Presupuestos mensuales por categoría \\
  Ingresos                   & 1--500    & Registro de ingresos del usuario \\
  Gastos                     & 0--5000   & Gastos manuales y generados por OCR \\
  Tickets                    & 0--1000   & Tickets escaneados con metadatos OCR \\
  Productos\_Tickets         & 0--10000  & Ítems extraídos de cada ticket \\
  Deudas                     & 0--20     & Deudas con seguimiento de pagos \\
  Alertas                    & 0--200    & Notificaciones automáticas del sistema \\
  Configuracion\_Notif.      & 1         & Preferencias de notificaciones por usuario \\
  Configuracion\_Inicial     & 1         & Configuración del onboarding inicial \\
  \bottomrule
\end{longtable}

\section{Descripción de Tablas}

\subsection{Tabla: Usuarios}

\begin{lstlisting}[style=sql, caption={DDL de la tabla Usuarios}, label={lst:tbl-usuarios}]
CREATE TABLE Usuarios (
  id_usuario            INT          PRIMARY KEY AUTO_INCREMENT,
  nombre                VARCHAR(100) NOT NULL,
  email                 VARCHAR(100) UNIQUE NOT NULL,
  password_hash         VARCHAR(255) NOT NULL,
  onboarding_completado BOOLEAN      DEFAULT FALSE,
  setup_completado      BOOLEAN      DEFAULT FALSE,
  fecha_creacion        TIMESTAMP    DEFAULT CURRENT_TIMESTAMP
) ENGINE=InnoDB DEFAULT CHARSET=utf8mb4;
\end{lstlisting}

\begin{longtable}{>{\ttfamily}p{4cm} >{\ttfamily}p{2.5cm} p{7cm}}
  \caption{Columnas de la tabla Usuarios}
  \label{tab:cols-usuarios} \\
  \toprule
  \textbf{Columna} & \textbf{Tipo} & \textbf{Descripción} \\
  \midrule
  id\_usuario            & INT PK AI  & Identificador único autoincrementable \\
  nombre                 & VARCHAR(100) & Nombre completo del usuario \\
  email                  & VARCHAR(100) UNIQUE & Email de acceso; índice único \\
  password\_hash         & VARCHAR(255) & Hash bcrypt de la contraseña (10 rounds) \\
  onboarding\_completado & BOOLEAN & \texttt{TRUE} cuando el usuario completó el flujo de bienvenida \\
  setup\_completado      & BOOLEAN & \texttt{TRUE} cuando guardó la configuración inicial \\
  fecha\_creacion        & TIMESTAMP & Fecha y hora de registro \\
  \bottomrule
\end{longtable}

\subsection{Tabla: Categorias}

\begin{lstlisting}[style=sql, caption={DDL de la tabla Categorias}]
CREATE TABLE Categorias (
  id_categoria    INT          PRIMARY KEY AUTO_INCREMENT,
  id_usuario      INT          NOT NULL,
  nombre_categoria VARCHAR(50) NOT NULL,
  tipo_categoria   ENUM('predefinida','personalizada') DEFAULT 'personalizada',
  icono           VARCHAR(50)  DEFAULT 'category',
  color           VARCHAR(7)   DEFAULT '#6B7280',
  FOREIGN KEY (id_usuario) REFERENCES Usuarios(id_usuario) ON DELETE CASCADE
) ENGINE=InnoDB DEFAULT CHARSET=utf8mb4;
\end{lstlisting}

Las 10 categorías predefinidas creadas automáticamente al registrar un usuario son: Alimentos, Transporte, Salud, Entretenimiento, Servicios, Ropa, Hogar, Educación, Restaurantes y Otros.

\subsection{Tabla: Presupuestos}

\begin{lstlisting}[style=sql, caption={DDL de la tabla Presupuestos}]
CREATE TABLE Presupuestos (
  id_presupuesto  INT           PRIMARY KEY AUTO_INCREMENT,
  id_usuario      INT           NOT NULL,
  id_categoria    INT           NOT NULL,
  monto_asignado  DECIMAL(10,2) NOT NULL,
  monto_utilizado DECIMAL(10,2) DEFAULT 0.00,
  periodo         ENUM('mensual','semanal','anual') DEFAULT 'mensual',
  fecha_inicio    DATE          NOT NULL,
  fecha_fin       DATE          NOT NULL,
  activo          BOOLEAN       DEFAULT TRUE,
  FOREIGN KEY (id_usuario)   REFERENCES Usuarios(id_usuario)   ON DELETE CASCADE,
  FOREIGN KEY (id_categoria) REFERENCES Categorias(id_categoria)
) ENGINE=InnoDB DEFAULT CHARSET=utf8mb4;
\end{lstlisting}

\subsection{Tabla: Ingresos}

\begin{lstlisting}[style=sql, caption={DDL de la tabla Ingresos}]
CREATE TABLE Ingresos (
  id_ingreso      INT           PRIMARY KEY AUTO_INCREMENT,
  id_usuario      INT           NOT NULL,
  nombre_ingreso  VARCHAR(100)  NOT NULL,
  monto           DECIMAL(10,2) NOT NULL,
  es_recurrente   BOOLEAN       DEFAULT FALSE,
  fecha           DATE          NOT NULL,
  fecha_creacion  TIMESTAMP     DEFAULT CURRENT_TIMESTAMP,
  FOREIGN KEY (id_usuario) REFERENCES Usuarios(id_usuario) ON DELETE CASCADE
) ENGINE=InnoDB DEFAULT CHARSET=utf8mb4;
\end{lstlisting}

\subsection{Tabla: Gastos}

\begin{lstlisting}[style=sql, caption={DDL de la tabla Gastos}]
CREATE TABLE Gastos (
  id_gasto        INT           PRIMARY KEY AUTO_INCREMENT,
  id_usuario      INT           NOT NULL,
  id_ticket       INT           NULL,    -- NULL si es gasto manual
  id_categoria    INT           NOT NULL,
  nombre_gasto    VARCHAR(100)  NOT NULL,
  monto           DECIMAL(10,2) NOT NULL,
  fecha           DATE          NOT NULL,
  extraido_por_ocr BOOLEAN      DEFAULT FALSE,
  notas           TEXT          NULL,
  fecha_creacion  TIMESTAMP     DEFAULT CURRENT_TIMESTAMP,
  FOREIGN KEY (id_usuario)   REFERENCES Usuarios(id_usuario)  ON DELETE CASCADE,
  FOREIGN KEY (id_ticket)    REFERENCES Tickets(id_ticket)     ON DELETE SET NULL,
  FOREIGN KEY (id_categoria) REFERENCES Categorias(id_categoria)
) ENGINE=InnoDB DEFAULT CHARSET=utf8mb4;
\end{lstlisting}

\subsection{Tabla: Tickets}

\begin{lstlisting}[style=sql, caption={DDL de la tabla Tickets}]
CREATE TABLE Tickets (
  id_ticket       INT           PRIMARY KEY AUTO_INCREMENT,
  id_usuario      INT           NOT NULL,
  nombre_tienda   VARCHAR(100)  DEFAULT 'Desconocida',
  total           DECIMAL(10,2) DEFAULT 0.00,
  metodo_entrada  ENUM('ocr','manual') DEFAULT 'ocr',
  confianza_ocr   DECIMAL(4,3)  NULL,    -- 0.000 a 1.000
  rutas_imagenes  JSON          NULL,    -- arreglo de rutas de archivo
  fecha_compra    DATE          NOT NULL,
  fecha_creacion  TIMESTAMP     DEFAULT CURRENT_TIMESTAMP,
  FOREIGN KEY (id_usuario) REFERENCES Usuarios(id_usuario) ON DELETE CASCADE
) ENGINE=InnoDB DEFAULT CHARSET=utf8mb4;
\end{lstlisting}

\subsection{Tabla: Productos\_Tickets}

\begin{lstlisting}[style=sql, caption={DDL de la tabla Productos\_Tickets}]
CREATE TABLE Productos_Tickets (
  id_producto     INT           PRIMARY KEY AUTO_INCREMENT,
  id_ticket       INT           NOT NULL,
  id_categoria    INT           NULL,
  nombre_producto VARCHAR(150)  NOT NULL,
  precio          DECIMAL(10,2) NOT NULL,
  FOREIGN KEY (id_ticket)    REFERENCES Tickets(id_ticket)       ON DELETE CASCADE,
  FOREIGN KEY (id_categoria) REFERENCES Categorias(id_categoria) ON DELETE SET NULL
) ENGINE=InnoDB DEFAULT CHARSET=utf8mb4;
\end{lstlisting}

\subsection{Tabla: Deudas}

\begin{lstlisting}[style=sql, caption={DDL de la tabla Deudas}]
CREATE TABLE Deudas (
  id_deuda             INT           PRIMARY KEY AUTO_INCREMENT,
  id_usuario           INT           NOT NULL,
  nombre_deuda         VARCHAR(100)  NOT NULL,
  monto_total          DECIMAL(10,2) NOT NULL,
  monto_pagado         DECIMAL(10,2) DEFAULT 0.00,
  pago_mensual         DECIMAL(10,2) NOT NULL,
  duracion_meses       INT           NOT NULL,
  fecha_siguiente_pago DATE          NOT NULL,
  activa               BOOLEAN       DEFAULT TRUE,
  fecha_creacion       TIMESTAMP     DEFAULT CURRENT_TIMESTAMP,
  FOREIGN KEY (id_usuario) REFERENCES Usuarios(id_usuario) ON DELETE CASCADE
) ENGINE=InnoDB DEFAULT CHARSET=utf8mb4;
\end{lstlisting}

\subsection{Tabla: Alertas}

\begin{lstlisting}[style=sql, caption={DDL de la tabla Alertas}]
CREATE TABLE Alertas (
  id_alerta   INT          PRIMARY KEY AUTO_INCREMENT,
  id_usuario  INT          NOT NULL,
  tipo_alerta VARCHAR(50)  NOT NULL,  -- 'presupuesto', 'deuda', 'gestion'
  titulo      VARCHAR(100) NOT NULL,
  mensaje     TEXT         NOT NULL,
  leida       BOOLEAN      DEFAULT FALSE,
  fecha       TIMESTAMP    DEFAULT CURRENT_TIMESTAMP,
  FOREIGN KEY (id_usuario) REFERENCES Usuarios(id_usuario) ON DELETE CASCADE
) ENGINE=InnoDB DEFAULT CHARSET=utf8mb4;
\end{lstlisting}

\subsection{Tabla: Configuracion\_Notificaciones}

\begin{lstlisting}[style=sql, caption={DDL de Configuracion\_Notificaciones}]
CREATE TABLE Configuracion_Notificaciones (
  id_config            INT     PRIMARY KEY AUTO_INCREMENT,
  id_usuario           INT     UNIQUE NOT NULL,
  presupuesto_excedido BOOLEAN DEFAULT TRUE,
  recordatorio_gestion BOOLEAN DEFAULT TRUE,
  pago_deuda_proximo   BOOLEAN DEFAULT TRUE,
  FOREIGN KEY (id_usuario) REFERENCES Usuarios(id_usuario) ON DELETE CASCADE
) ENGINE=InnoDB DEFAULT CHARSET=utf8mb4;
\end{lstlisting}

\subsection{Tabla: Configuracion\_Inicial}

\begin{lstlisting}[style=sql, caption={DDL de Configuracion\_Inicial}]
CREATE TABLE Configuracion_Inicial (
  id_config               INT          PRIMARY KEY AUTO_INCREMENT,
  id_usuario              INT          UNIQUE NOT NULL,
  modo_administracion     ENUM('unico','multiple') DEFAULT 'unico',
  distribucion_categorias JSON         NULL,
  ingreso_mensual_inicial DECIMAL(10,2) NULL,
  FOREIGN KEY (id_usuario) REFERENCES Usuarios(id_usuario) ON DELETE CASCADE
) ENGINE=InnoDB DEFAULT CHARSET=utf8mb4;
\end{lstlisting}

\section{Triggers Automáticos}

El sistema utiliza dos triggers MySQL que implementan lógica de negocio a nivel de base de datos, garantizando la integridad incluso si se insertan datos por fuera de la API.

\subsection{Trigger: actualizar\_presupuesto\_gasto}

\begin{lstlisting}[style=sql, caption={Trigger de actualización de presupuesto al registrar gasto}, label={lst:trigger1}]
DELIMITER $$
CREATE TRIGGER actualizar_presupuesto_gasto
AFTER INSERT ON Gastos
FOR EACH ROW
BEGIN
  UPDATE Presupuestos
  SET monto_utilizado = monto_utilizado + NEW.monto
  WHERE id_usuario   = NEW.id_usuario
    AND id_categoria = NEW.id_categoria
    AND activo       = TRUE
    AND fecha_inicio <= NEW.fecha
    AND fecha_fin    >= NEW.fecha;
END$$
DELIMITER ;
\end{lstlisting}

\subsection{Trigger: verificar\_presupuesto\_excedido}

\begin{lstlisting}[style=sql, caption={Trigger de alerta cuando el presupuesto alcanza el 90\%}, label={lst:trigger2}]
DELIMITER $$
CREATE TRIGGER verificar_presupuesto_excedido
AFTER UPDATE ON Presupuestos
FOR EACH ROW
BEGIN
  IF NEW.monto_utilizado >= (NEW.monto_asignado * 0.90) THEN
    INSERT INTO Alertas (id_usuario, tipo_alerta, titulo, mensaje)
    VALUES (
      NEW.id_usuario,
      'presupuesto',
      'Presupuesto al limite',
      CONCAT('Tu presupuesto ha alcanzado el ',
             ROUND((NEW.monto_utilizado/NEW.monto_asignado)*100, 0),
             '% del limite asignado.')
    );
  END IF;
END$$
DELIMITER ;
\end{lstlisting}

\section{Diagrama Entidad-Relación}

\begin{figure}[H]
  \centering
  \begin{tikzpicture}[
    entity/.style={rectangle, draw, fill=blue!8, rounded corners=2pt,
                   minimum width=3cm, minimum height=0.8cm, font=\small\bfseries},
    rel/.style={-Stealth, thick, gray},
    label/.style={font=\tiny, midway, fill=white, inner sep=1pt}
  ]
    \node[entity] (usr)  at (0,0)     {Usuarios};
    \node[entity] (cat)  at (-4,-2.5) {Categorias};
    \node[entity] (pre)  at (-2,-5)   {Presupuestos};
    \node[entity] (ing)  at (0,-2.5)  {Ingresos};
    \node[entity] (gas)  at (2,-5)    {Gastos};
    \node[entity] (tkt)  at (6,-2.5)  {Tickets};
    \node[entity] (pro)  at (6,-5)    {Productos\_Tickets};
    \node[entity] (deu)  at (4,-2.5)  {Deudas};
    \node[entity] (ale)  at (-4,0)    {Alertas};
    \node[entity] (not)  at (4,0)     {Config\_Notif.};
    \node[entity] (ini)  at (0,2)     {Config\_Inicial};

    \draw[rel] (usr) -- node[label,left]{1:N} (cat);
    \draw[rel] (usr) -- node[label,right]{1:N} (ing);
    \draw[rel] (usr) -- node[label]{1:N} (deu);
    \draw[rel] (usr) -- node[label,left]{1:N} (ale);
    \draw[rel] (usr) -- node[label,right]{1:1} (not);
    \draw[rel] (usr) -- node[label,right]{1:1} (ini);
    \draw[rel] (usr) -- node[label,left]{1:N} (tkt);
    \draw[rel] (cat) -- node[label]{1:N} (pre);
    \draw[rel] (cat) -- node[label,right]{1:N} (gas);
    \draw[rel] (tkt) -- node[label,right]{1:N} (pro);
    \draw[rel] (tkt) -- node[label]{1:1} (gas);
    \draw[rel] (cat) -- node[label]{1:N} (pro);
  \end{tikzpicture}
  \caption{Diagrama Entidad-Relación simplificado de finara\_db}
  \label{fig:er-diagram}
\end{figure}
